\documentclass[12pt]{article}
\usepackage[T1]{fontenc}
\usepackage[utf8]{inputenc}
\usepackage{lmodern}
\usepackage{textcomp}
\usepackage{enumitem}
\usepackage{geometry}
\geometry{margin=1in}
\begin{document}
\begin{center}
Answer Key - Economics - K-12 Level Assessment 15 \
\small (Answers are ordered exactly as in the question paper.)
\end{center}

\section*{Multiple Choice}
\begin{enumerate}[label=\arabic*.]
\item \textbf{Answer:} D
\\ \textit{Options:} A) direct relation between unemployment and inflation; B) direct relation between price and quantity demanded; C) inverse relation between price and quantity demanded; D) inverse relation between unemployment and inflation
\\ \textit{Note:} The short-run Phillips curve illustrates that as unemployment decreases, inflation tends to increase, reflecting an inverse relationship between the two economic indicators due to the trade-off between inflation and unemployment in the short run.
\item \textbf{Answer:} D
\\ \textit{Options:} A) The price floor shifts the demand curve to the left.; B) An effective floor creates a shortage of the good.; C) The price floor shifts the supply curve of the good to the right.; D) To be an effective floor, it must be set above the equilibrium price.
\\ \textit{Note:} A price floor is effective only when set above the equilibrium price because it prevents the price from falling to a level where supply meets demand, ensuring that producers receive a higher price for their goods.
\item \textbf{Answer:} D
\\ \textit{Options:} A) the price level is constant with an increase in aggregate demand.; B) the price level falls with an increase in aggregate supply.; C) the price level is constant with an increase in long-run aggregate supply.; D) the price level rises with an increase in aggregate demand.
\\ \textit{Note:} The correct answer is based on the fact that when the price level rises due to an increase in aggregate demand, it reduces the real purchasing power of consumers, thereby diminishing the overall impact of the spending multiplier on the economy.
\item \textbf{Answer:} A
\\ \textit{Options:} A) the United States to be a net exporter of copper.; B) the United States to impose a tariff on imported copper to protect domestic producers.; C) the demand for U.S. copper to fall.; D) a growing trade deficit in the United States in goods and services.
\\ \textit{Note:} When the world price of copper is higher than the domestic price, it incentivizes U.S. producers to export copper to take advantage of the higher prices abroad, leading to the United States becoming a net exporter of copper.
\item \textbf{Answer:} D
\\ \textit{Options:} A) is thought to slightly overestimate the inflation rate; B) uses base year quantities in its calculations; C) incorporates both current year prices and base year prices; D) incorporates current year quantities in its calculations
\\ \textit{Note:} The GDP Deflator measures the price level of all new, domestically produced goods and services in an economy by using current year quantities, while the CPI calculates inflation based on a fixed basket of goods and services from a base year, making the former more reflective of current economic conditions.
\end{enumerate}

\section*{True / False}
\begin{enumerate}[label=\arabic*.]
\setcounter{enumi}{5}
\item \textbf{Answer:} False
\\ \textit{Options:} A) True; B) False
\\ \textit{Note:} Ceteris paribus, which translates to "all other things being equal," means that when analyzing the effect of one variable, all other variables are held constant, indicating that not everything can change at any time.
\item \textbf{Answer:} True
\\ \textit{Options:} A) True; B) False
\\ \textit{Note:} Air is essential for survival, making it highly valuable in total utility, yet its abundance leads to a very low marginal utility, as additional units provide little extra benefit.
\item \textbf{Answer:} False
\\ \textit{Options:} A) True; B) False
\\ \textit{Note:} The Federal Reserve primarily uses monetary policy tools, such as interest rates and open market operations, rather than tariffs and quotas, which are trade policies typically set by the government.
\item \textbf{Answer:} False
\\ \textit{Options:} A) True; B) False
\\ \textit{Note:} An increase in aggregate demand typically leads to lower unemployment in the short run but can also result in higher inflation rates as demand outpaces supply, making the statement false.
\item \textbf{Answer:} False
\\ \textit{Options:} A) True; B) False
\\ \textit{Note:} The statement is false because national income is calculated by considering the total value of goods and services produced in a country, which includes factors of production such as labor and capital, not just household income.
\end{enumerate}

\section*{Long Form Response}
\begin{enumerate}[label=\arabic*.]
\setcounter{enumi}{10}
\item \textbf{Answer:} In the circular flow model of a private closed economy, firms and households play crucial roles in the exchange of goods and revenues. Firms produce goods and services that are supplied to households, which are the consumers in the economy. In return for these goods, households provide firms with revenue through their purchases. This exchange creates a continuous flow where firms rely on household spending to generate income, while households depend on firms for the goods and services they need. This interdependence illustrates how both sectors work together to sustain the economy.
\\ \textit{Note:} The answer correctly explains the interdependent relationship between firms and households in the circular flow model, emphasizing that firms supply goods and services to households, while households generate revenue for firms through their consumption.
\item \textbf{Answer:} The introduction of a new production technique that lowers costs can significantly influence economic growth by increasing efficiency and productivity. For example, if a factory adopts a new automated system that reduces labor costs, it can produce goods faster and cheaper. This not only allows the factory to increase its output but also enables it to lower prices, making its products more competitive in the market. As a result, more consumers may buy these products, which can lead to higher sales, more jobs, and ultimately, a boost in the overall economy. This cycle of increased production and consumption can drive economic growth and improve living standards.
\\ \textit{Note:} correct because it clearly illustrates how a new production technique can enhance efficiency and productivity, leading to lower costs, competitive pricing, increased consumer demand, and ultimately contributing to economic growth.
\end{enumerate}

\vfill
\noindent\textit{End of Answer Key}
\end{document}